\section{Discussion and Future Work}
\label{sec:conclusion}

OpenSkynet demonstrates that self-healing, skill learning, and tiered memory can be integrated into a unified browser automation framework. The system's architecture addresses three persistent limitations of LLM-powered agents: execution brittleness, lack of cross-session learning, and single-agent context window constraints.

\subsection{Limitations}

Several limitations motivate future work. First, the reliance on LLMs for both execution and healing creates a circular dependency: the same model that failed to execute the original task is asked to diagnose and repair the failure. Exploring external validation mechanisms---such as execution-based verifiers or programmatic constraint checkers---could reduce this dependency. Second, the skill learning pipeline currently generates text-based steps; incorporating structured execution traces with element selectors, wait conditions, and error handlers could improve robustness. Third, the memory system's importance scoring is heuristic-based; adopting learned importance models (e.g., fine-tuned classifiers or RL-trained critics) could improve promotion accuracy.

\subsection{Future Directions}

\begin{itemize}[leftmargin=*,itemsep=2pt]
    \item \textbf{Full-scale benchmark evaluation.} Running the Evil Skill Generalization Benchmark across all 60 tasks with multiple LLM providers will produce quantitative results on self-healing efficacy and skill generalization rates.
    \item \textbf{Standard benchmark integration.} Evaluating OpenSkynet on WebArena and VisualWebArena will provide apples-to-apples comparison with existing browser agent systems.
    \item \textbf{Ablation studies.} Systematic ablation of individual components (healing on/off, skill learning on/off, memory tier on/off) will isolate the contribution of each mechanism.
    \item \textbf{RL fine-tuning.} Combining the framework's trajectory database with reinforcement learning (inspired by WebRL's self-evolving curriculum) could train specialized models that are inherently more robust to page changes.
    \item \textbf{Skill marketplace.} OpenSkynet's skills hub currently supports 470+ community skills. Formalizing skill sharing with trust levels, verification protocols, and reputation systems could create a collaborative ecosystem for browser automation skills.
    \item \textbf{Cross-modal skill transfer.} Extending the Trace-to-Skill module to accept video demonstrations (rather than just browser recordings) would enable users to ``teach'' the system by recording their screen, further lowering the automation barrier.
\end{itemize}

\subsection{Conclusion}

We presented OpenSkynet, a framework that combines self-healing skill execution, automated skill learning, tiered memory, and multi-agent orchestration for persistent browser automation. The system's architecture is designed to improve over time: skills become more robust through healing, memory accumulates cross-session knowledge, and subagents enable scalable task decomposition. The proposed Evil Skill Generalization Benchmark provides a principled testbed for evaluating self-healing and generalization in adversarial web environments. We release OpenSkynet as open-source software to facilitate further research and practical adoption in the browser automation community.
